\begin{abstract}
Human oversight of machine learning systems is often treated as a simple fallback: when model confidence is low, defer to a human. In financial fraud review, however, such strategies can be inadequate because wrong-but-confident model failures can be especially costly, fraud patterns are heterogeneous, and oversight resources are limited and uneven. We present \textbf{MAGD-Fraud} (Multi-evidence Assurance-Guided Deferral), a routing framework for financial fraud alert review that combines multiple deployable assurance signals rather than relying on distance uncertainty or score confidence alone. MAGD-Fraud integrates numerical confidence, calibration risk, embedding-based distance uncertainty, local historical neighbourhood reliability, and wrong-confident AI risk into a composite assurance score, and uses this score together with adaptive thresholding and expert-aware routing to select among AI, a synthetic expert, or escalation. The framework further supports learned and constrained policy variants that optimize validation-time routing weights under overreliance, capacity, fairness, and audit-oriented objectives. We implement MAGD-Fraud in the FiFAR benchmark setting and evaluate it against AI-only, threshold-based deferral, expert baselines, and a learning-to-defer baseline. Beyond standard predictive metrics, we report cost-sensitive loss, overreliance, correct rejection, wrong-confident avoidance, and audit coverage. We position MAGD-Fraud as a computational Human--AI assurance prototype for fraud review: the expert panels are synthetic experts and the dashboard is a prototype interface, so the results should not be interpreted as a human-subject validation study.
\end{abstract}

\section{Introduction}
\label{sec:introduction}

Machine learning systems are increasingly used to prioritize and support high-stakes review tasks, including financial fraud alert investigation. In such settings, an automated model is rarely the only decision maker. Instead, institutions often rely on a combination of model predictions, human reviewers, escalation processes, and operational policies. This creates a central Human--AI question: \emph{when should the system rely on AI, when should it defer to a reviewer, and when should it escalate?}

A common answer is to use model confidence or uncertainty as a deferral trigger. If the model appears uncertain, defer; if it appears confident, rely on the model. While intuitive, this approach can be fragile in fraud review. Fraud cases are often rare, heterogeneous, and locally irregular. A case may appear geometrically familiar under one uncertainty measure but still resemble a neighborhood of historically misclassified cases. A model may also be globally well calibrated while remaining wrong and highly confident on precisely the types of cases for which human oversight matters most. In short, \emph{single-signal uncertainty is often too weak a basis for high-stakes routing}.

This paper proposes \textbf{MAGD-Fraud} (Multi-evidence Assurance-Guided Deferral), a Human--AI routing framework that extends uncertainty-based deferral into a broader assurance-guided decision layer. Rather than treating distance uncertainty or score confidence as sufficient, MAGD-Fraud combines multiple deployable assurance signals: numerical confidence, calibration risk, embedding-based distance uncertainty, local historical neighbourhood reliability, and wrong-confident AI risk. These signals are aggregated into a case-level assurance score, which is then used together with adaptive thresholding and expert-aware expected-cost routing to choose among three routes: use AI, defer to a selected synthetic expert, or escalate to a panel of reliable experts.

Our motivation is both technical and socio-technical. Technically, we seek a routing policy that is better matched to the structure of fraud data than centroid-only uncertainty. Socio-technically, we seek a policy that is evaluated not only by predictive quality, but also by whether it reduces harmful overreliance on wrong AI predictions, supports justified expert intervention, and preserves auditability. In this sense, MAGD-Fraud is less a new classifier than a new \emph{assurance architecture} layered on top of a fraud predictor.

We make three main contributions.

\paragraph{First,} we introduce a \emph{multi-evidence assurance representation} for fraud-review routing. The representation combines score-based, geometric, and local historical signals, and explicitly models wrong-confident AI risk rather than assuming confidence is inherently trustworthy.

\paragraph{Second,} we formulate routing as an \emph{assurance-guided deferral policy} rather than a simple thresholding rule. MAGD-Fraud supports heuristic, learned, and constrained policy variants, and its constrained variant incorporates penalties for overreliance, capacity violations, fairness penalties, and audit gaps.

\paragraph{Third,} we evaluate the framework using metrics that are more appropriate for Human--AI assurance than accuracy alone. In addition to predictive performance, we report cost-sensitive loss, overreliance, correct rejection, wrong-confident avoidance, coverage, escalation behavior, and audit coverage. We also generate claim--evidence and audit artifacts to support traceable analysis of routing behavior.

Our work is motivated by the need for more realistic computational studies of Human--AI oversight in fraud review, but it is important to be clear about scope. The FiFAR environment used here contains \emph{synthetic experts}, not real human reviewers. The user-facing dashboard in the repository is a prototype interface, not a validated human decision-support tool. Accordingly, we frame MAGD-Fraud as a computational research prototype for Human--AI assurance, not as evidence of human-subject efficacy or deployment readiness.

The rest of the paper proceeds as follows. Section~\ref{sec:related_work} situates MAGD-Fraud relative to uncertainty estimation, selective prediction, learning to defer, and Human--AI decision-making research. Section~\ref{sec:method} presents the MAGD-Fraud architecture, including the assurance signals, routing policy, and learned and constrained variants. Section~\ref{sec:experimental_setup} describes the FiFAR benchmark setting, baselines, metrics, and leakage boundaries. We then report empirical results, ablations, statistical tests, and audit-oriented analyses.
